Residual access is restricted to primary survivors, reducing coordinate
work from $O(n_cd)$ to $O(cd)$. Lane cooperation accelerates that work
at a fixed budget.

\subsection{Reconstructed scoring and the query-only constant}

Let $\widehat p(x)$ be the primary reconstruction and $\widehat r(x)$
the decoded residual in transformed coordinates. The encoder stores
$\gamma(x)=2\langle\widehat p(x),\widehat r(x)\rangle$.
The implemented refined ranking score is
\begin{align}
 \refineddist(q,x)
 &=\primarydist(q,x)+\|q'-\widehat r(x)\|_2^2+\gamma(x)
 \label{eq:5-refined-score}\\
 &=\|q'-\widehat p(x)-\widehat r(x)\|_2^2+\|q'\|_2^2.
 \label{eq:5-refined-constant}
\end{align}
Expanding the norms proves the second equality. The returned score is
reconstructed squared distance plus a query-only constant, which
preserves ordering while retaining reconstruction error \cite{jhq}.

Residual coordinates use scalar codebook entries shared within a
subspace. The scalar path walks each candidate's dimensions and sums
their squared differences from $q'$. Query coordinates and, when space
permits, the residual codebook are staged in shared memory; the stored
correction and primary score are added afterwards.

\subsection{Mapping lanes to words and reducing partial sums}

For $\lrefine=L\in\{1,8,16,32\}$, a scan-independent refinement block
assigns $L$ adjacent lanes to a candidate. A 32-bit residual word holds
$d_{\rm word}=32/B_R$ coordinate codes, where $B_R$ is the residual bit width.
The tested eight-bit path has $d_{\rm word}=4$; the four-bit path has $d_{\rm word}=8$.
Lane $a$ reads word indices $a,a+L,a+2L,\ldots$, decodes their
coordinates and accumulates a partial sum. A butterfly shuffle reduction
combines the $L$ partials, and lane zero writes the refined score.

\begin{figure}[tbp]
\centering
% Lane mapping for L=16; schematic, not a measured trace.
\begin{tikzpicture}[font=\small,
 box/.style={draw=ThesisNavy!65,rounded corners=2pt,align=center,
   text width=5.8cm,minimum height=1.5cm,inner sep=5pt},
 arr/.style={-{Latex[length=2mm]},draw=ThesisNavy!80,thick}]
\node[box,fill=ThesisNavy!5] (a) at (0,0)
 {Warp lanes 0--15: candidate $x_0$\\
  subgroup lane $a$: words $a,a+16,\ldots$};
\node[box,fill=ThesisNavy!5] (b) at (7,0)
 {Warp lanes 16--31: candidate $x_1$\\
  subgroup lane $a$: words $a,a+16,\ldots$};
\node[box] (ar) at (0,-2.2)
 {Decode coordinates; accumulate partials\\
  Shuffle offsets $8,4,2,1$};
\node[box] (br) at (7,-2.2)
 {Decode coordinates; accumulate partials\\
  Shuffle offsets $8,4,2,1$};
\node[box,text width=11.5cm,minimum height=1.15cm] (out) at (3.5,-4.25)
 {Lane zero of each subgroup writes\\
  $\refineddist=\primarydist+\text{residual squared-difference sum}+\gamma(x)$};
\draw[arr] (a) -- (ar);
\draw[arr] (b) -- (br);
\draw[arr] (ar.south) -- (out.north west);
\draw[arr] (br.south) -- (out.north east);
\end{tikzpicture}

\caption[Sixteen-lane residual refinement]{Sixteen-lane residual refinement. A warp handles two candidates;
lane $a$ within a candidate reads residual words $a,a+16,\ldots$.
With eight-bit residual codes, each 32-bit word decodes four coordinates.
Subgroup reductions combine
partial scores before adding the retained primary score and correction.
Candidate positions can be scattered even when each subgroup's word
access is contiguous.}
\label{fig:5-lane-refinement}
\end{figure}

The launch requires a word-aligned residual row and $D_s$ divisible by
$d_{\rm word}$, so a word never crosses a subspace's codebook boundary. All lanes
execute the shuffle sequence, including lanes assigned to invalid or
out-of-range survivor slots; these contribute zero before the output
is marked infinite. This avoids making warp participation depend on
whether a particular candidate exists. At $L=16$ a warp handles two
candidates, rather than one candidate per lane.

The experiment separates two changes. L0 is the original scalar
byte-load path; L1 uses word loads with one lane per candidate, preserving
its accumulation order. L8/L16/L32 additionally distribute coordinate
work across lanes. This shortens each lane's serial walk and provides
useful work when $c$ is small, but adds shuffles and reduces the number
of independent candidates handled concurrently. Larger $L$ is therefore
not automatically better. Irregular survivor positions and shared/global
codebook placement remain relevant costs.

\subsection{Final selection and the numerical boundary}

For at least $k$ valid survivors with finite refined scores and unique
database IDs, the final kernel performs exact top-$k$ selection under
$(\refineddist,\mathrm{database\ ID})$. One block per query loads the
$c$ scores and repeatedly reduces the minimum, emits its ID and score,
and removes that score before the next round. All $c$ slots are scanned
at each round; there is no fixed per-thread candidate quota. In a
degenerate pool, selected scores become infinite but positions remain
valid: after finite entries run out, the smallest valid ID can repeat
with infinity. An entirely invalid pool returns invalid IDs and infinity.
The evaluation satisfies the finite-entry contract.

The final ID tie rule differs from the physical-position rule used for
primary retention. An exact final selector cannot recover candidates
discarded by routing or screening, nor remove quantisation error. Wider
lane reductions also change floating-point association before selection,
which can move near-tied candidates even when the real-arithmetic
formula is unchanged.

\Cref{sec:7-3} reports the load-only and cooperative controls,
numerical differences, kernel times and end-to-end gains separately.
